% ===========================================================================
%  Devoir surveillé nº 4 — fin de la partie Traitement des données
%  Sujet complet, au format de l'épreuve
% ===========================================================================
\chapter*{Devoir surveillé nº\,4 — sujet complet}
\markboth{Devoir surveillé nº\,4}{Devoir surveillé nº\,4}
\addcontentsline{toc}{chapter}{Devoir surveillé nº\,4 — sujet complet}

\begin{devoirsurveille}[portée : tout le programme]
\textbf{Durée : 2 h 15.} Les deux exercices et le problème sont obligatoires.
Calculatrice scientifique non programmable autorisée. Les résultats des
probabilités seront donnés sous forme de fraction irréductible.

\vspace{0.7em}
\noindent\textbf{\sffamily EXERCICE 1} \hfill \textup{(05 points)}

Une boîte contient $10$ jetons indiscernables au toucher : $4$ jetons verts,
$3$ jetons jaunes et $3$ jetons rouges.

\begin{enumerate}[label=\textbf{\arabic*.}, leftmargin=*, itemsep=0.25em]
  \item On tire au hasard et simultanément $3$ jetons de la boîte.
    \begin{enumerate}[label=\textbf{\alph*)}, leftmargin=1.4em, itemsep=0.15em]
      \item Déterminer le nombre de cas possibles. \bareme{0,5}
      \item Calculer la probabilité de l'événement
            $A$ : « obtenir trois jetons de même couleur ». \bareme{1}
      \item Calculer la probabilité de l'événement
            $B$ : « obtenir exactement deux jetons verts ». \bareme{1}
      \item Calculer la probabilité de l'événement
            $C$ : « obtenir au moins un jeton rouge ». \bareme{1}
    \end{enumerate}
  \item On remet les jetons dans la boîte, puis on tire successivement et
        avec remise $3$ jetons. Calculer la probabilité de l'événement
        $D$ : « obtenir, dans l'ordre, un jeton vert, un jaune et un rouge ».
        \bareme{1,5}
\end{enumerate}

\vspace{0.7em}
\noindent\textbf{\sffamily EXERCICE 2} \hfill \textup{(05 points)}

On considère la suite $\left(U_{n}\right)_{n \in \N}$ définie par
$U_{0} = 1$ et, pour tout entier naturel $n$,
\[
  U_{n+1} = \frac12 U_{n}+3 .
\]

\begin{enumerate}[label=\textbf{\arabic*.}, leftmargin=*, itemsep=0.25em]
  \item Calculer $U_{1}$ et $U_{2}$. \bareme{0,5}
  \item On pose $V_{n} = U_{n}-6$ pour tout entier naturel $n$.
    \begin{enumerate}[label=\textbf{\alph*)}, leftmargin=1.4em, itemsep=0.15em]
      \item Montrer que $\left(V_{n}\right)$ est une suite géométrique dont
            on précisera la raison et le premier terme. \bareme{1,5}
      \item Exprimer $V_{n}$ puis $U_{n}$ en fonction de $n$.
            \bareme{1}
      \item Calculer $\limn U_{n}$. \bareme{0,5}
    \end{enumerate}
  \item Soit $S_{n} = V_{0}+V_{1}+\cdots+V_{n}$.
    \begin{enumerate}[label=\textbf{\alph*)}, leftmargin=1.4em, itemsep=0.15em]
      \item Exprimer $S_{n}$ en fonction de $n$. \bareme{1}
      \item En déduire $T_{n} = U_{0}+U_{1}+\cdots+U_{n}$. \bareme{0,5}
    \end{enumerate}
\end{enumerate}

\vspace{0.7em}
\noindent\textbf{\sffamily PROBLÈME} \hfill \textup{(10 points)}

On considère la fonction numérique $f$ définie sur
$\intervalleo{0}{+\infty}$ par
\[
  f(x) = \frac{1}{x}+\ln x ,
\]
et l'on note $(\mathcal C)$ sa courbe représentative dans un repère
orthonormé $\left(O\,;\,\veci\,,\,\vecj\right)$ d'unité graphique $2$ cm.

\begin{enumerate}[label=\textbf{\arabic*.}, leftmargin=*, itemsep=0.25em]
  \item \begin{enumerate}[label=\textbf{\alph*)}, leftmargin=1.4em, itemsep=0.15em]
      \item Calculer $\displaystyle\lim_{x \to 0^{+}} f(x)$ et en déduire
            une asymptote à $(\mathcal C)$. \bareme{0,75}
      \item Calculer $\limpinf f(x)$. \bareme{0,75}
    \end{enumerate}
  \item \begin{enumerate}[label=\textbf{\alph*)}, leftmargin=1.4em, itemsep=0.15em]
      \item Montrer que, pour tout $x>0$,
            $f'(x) = \dfrac{x-1}{x^{2}}$. \bareme{1,25}
      \item Étudier le signe de $f'(x)$ sur $\intervalleo{0}{+\infty}$.
            \bareme{0,75}
      \item Dresser le tableau de variation de $f$. \bareme{1,25}
    \end{enumerate}
  \item En déduire que, pour tout $x>0$, $f(x) \geq 1$. Que peut-on en
        conclure sur les points d'intersection de $(\mathcal C)$ avec l'axe
        des abscisses ? \bareme{1}
  \item Écrire une équation de la tangente $(T)$ à $(\mathcal C)$ au point
        d'abscisse $1$. \bareme{1}
  \item Calculer $f(0{,}5)$, $f(2)$ et $f(3)$ à $10^{-2}$ près, puis tracer
        $(T)$ et $(\mathcal C)$ sur $\intervalleof{0}{3}$. On donne
        $\ln 2 \approx 0{,}69$ et $\ln 3 \approx 1{,}10$.
        \bareme{1,75}
\end{enumerate}

\begin{enumerate}[label=\textbf{\arabic*.}, leftmargin=*, itemsep=0.25em, start=6]
  \item Soit $F$ la fonction définie sur $\intervalleo{0}{+\infty}$ par
        $F(x) = \ln x+x\ln x-x$.
    \begin{enumerate}[label=\textbf{\alph*)}, leftmargin=1.4em, itemsep=0.15em]
      \item Montrer que $F$ est une primitive de $f$ sur
            $\intervalleo{0}{+\infty}$. \bareme{0,75}
      \item Calculer, en $\text{cm}^{2}$, l'aire du domaine plan délimité
            par $(\mathcal C)$, l'axe des abscisses et les droites
            d'équations $x = 1$ et $x = e$. On donne $e \approx 2{,}72$.
            \bareme{1,25}
    \end{enumerate}
\end{enumerate}
\end{devoirsurveille}

\begin{corrige}[devoir surveillé nº 4 — exercice 1]
\textbf{1. a)} $\combi{3}{10} = \frac{10\times9\times8}{6} = 120$ cas
possibles.

\textbf{b)} Trois verts : $\combi{3}{4} = 4$ ; trois jaunes :
$\combi{3}{3} = 1$ ; trois rouges : $\combi{3}{3} = 1$. Total : $6$ cas,
donc $\proba{A} = \frac{6}{120} = \frac{1}{20}$.

\textbf{c)} Deux verts parmi quatre et un jeton parmi les six autres :
$\combi{2}{4}\times\combi{1}{6} = 6\times6 = 36$, donc
$\proba{B} = \frac{36}{120} = \frac{3}{10}$.

\textbf{d)} Le contraire est « aucun rouge », soit
$\combi{3}{7} = 35$ cas. Donc
$\proba{C} = 1-\frac{35}{120} = \frac{85}{120} = \frac{17}{24}$.

\textbf{2.} Cas possibles : $10^{3} = 1\,000$. Cas favorables :
$4\times3\times3 = 36$. Donc
$\proba{D} = \frac{36}{1\,000} = \frac{9}{250}$.
\end{corrige}

\begin{corrige}[devoir surveillé nº 4 — exercice 2]
\textbf{1.} $U_{1} = \frac12+3 = \frac72$ et
$U_{2} = \frac74+3 = \frac{19}{4}$.

\textbf{2. a)} $V_{n+1} = U_{n+1}-6 = \frac12U_{n}+3-6 = \frac12U_{n}-3
= \frac12\left(U_{n}-6\right) = \frac12 V_{n}$ : la suite est géométrique
de raison $\frac12$ et de premier terme $V_{0} = 1-6 = -5$.
\textbf{b)} $V_{n} = -5\left(\frac12\right)^{n}$ et
$U_{n} = 6-5\left(\frac12\right)^{n}$.
\textbf{c)} Comme $\left|\frac12\right|<1$, $\left(\frac12\right)^{n} \to 0$
et $U_{n} \to 6$.

\textbf{3. a)} $S_{n} = -5\times
\frac{1-\left(\frac12\right)^{n+1}}{1-\frac12}
= -10\left(1-\left(\frac12\right)^{n+1}\right)$.
\textbf{b)} Comme $U_{k} = V_{k}+6$ et qu'il y a $n+1$ termes :
$T_{n} = S_{n}+6(n+1)$.
\end{corrige}

\begin{corrige}[devoir surveillé nº 4 — problème]
\textbf{1. a)} En $0^{+}$, $\frac1x \to +\infty$ et $\ln x \to -\infty$ :
il faut lever l'indétermination. On écrit
\[
  f(x) = \frac{1+x\ln x}{x},
\]
dont le numérateur tend vers $1$ (car $x\ln x \to 0$) et le dénominateur
vers $0^{+}$ : donc $f(x) \to +\infty$. L'axe des ordonnées, d'équation
$x = 0$, est asymptote verticale à $(\mathcal C)$.
\textbf{b)} En $+\infty$, $\frac1x \to 0$ et $\ln x \to +\infty$ : donc
$f(x) \to +\infty$.

\textbf{2. a)} $f'(x) = -\frac{1}{x^{2}}+\frac1x
= \frac{-1+x}{x^{2}} = \frac{x-1}{x^{2}}$.
\textbf{b)} Le dénominateur $x^{2}$ est strictement positif : le signe de
$f'$ est celui de $x-1$, négatif sur $\intervalleo{0}{1}$, nul en $1$,
positif sur $\intervalleo{1}{+\infty}$.
\textbf{c)} $f$ décroît sur $\intervalleof{0}{1}$, croît sur
$\intervallefo{1}{+\infty}$ ; les deux limites valent $+\infty$ et le
minimum vaut $f(1) = 1+0 = 1$.

\textbf{3.} Le minimum de $f$ sur $\intervalleo{0}{+\infty}$ vaut $1$ :
donc $f(x) \geq 1 > 0$ pour tout $x>0$. La courbe reste entièrement
au-dessus de l'axe des abscisses : elle ne le coupe en aucun point.

\textbf{4.} $f(1) = 1$ et $f'(1) = 0$ : la tangente est horizontale,
$(T) : y = 1$.

\textbf{5.} $f(0{,}5) = 2-0{,}69 = 1{,}31$ ;
$f(2) = 0{,}5+0{,}69 = 1{,}19$ ; $f(3) \approx 0{,}33+1{,}10 = 1{,}43$. On
trace l'asymptote verticale, la tangente horizontale $y = 1$, le sommet
$(1\,;\,1)$, puis la courbe, en forme de vallée, qui descend depuis
l'infini, touche son minimum en $x = 1$, puis remonte.

\textbf{6. a)} $F'(x) = \frac1x+\left(\ln x+1\right)-1 = \frac1x+\ln x
= f(x)$. \checkmark
\textbf{b)} La fonction étant positive sur $\intervalle{1}{e}$,
\[
  \mathcal A = F(e)-F(1) = \left(1+e-e\right)-\left(0+0-1\right) = 2
  \text{ unités d'aire}.
\]
L'unité graphique valant $2$ cm, une unité d'aire vaut $4$ cm² :
$\mathcal A = 8$ cm².
\end{corrige}

\begin{notepourlaclasse}
Ce devoir est un sujet complet : il peut servir de baccalauréat blanc, tel
quel, au troisième trimestre. Sa difficulté est calibrée légèrement
au-dessus de la moyenne des sessions récentes, sur deux points précis — la
levée de l'indétermination en $0^{+}$ à la question 1.a) du problème, et la
question 3, qui demande de tirer une conséquence du tableau de variation
plutôt que de le dresser.

Si vous le donnez en cours d'année plutôt qu'en fin, remplacez la question
1.a) par « on admet que $\displaystyle\lim_{x \to 0^{+}} f(x) = +\infty$ » :
le reste du problème est alors accessible dès la fin du chapitre 9.

Barème indicatif de correction : comptez un demi-point, dans chaque
question, pour la \emph{phrase} de conclusion — asymptote, nature de la
suite, interprétation de l'aire. C'est la part que les élèves de la série
perdent le plus systématiquement, et la leur montrer chiffrée est le moyen
le plus efficace de la leur faire écrire.
\end{notepourlaclasse}
